% =====================================================================
%  MSACL DS301 Deep Learning · Quiz 5 · Convolutional Networks
%  Academic problem-set style (single-source, \ifsolution toggle)
%  SINGLE SOURCE — produces BOTH the student quiz and the answer key.
%
%  ANSWER TOGGLE
%  -------------
%  Student version (answers HIDDEN) — the default:
%       pdflatex quiz05_cnns.tex
%  Answer key (answers SHOWN):
%       pdflatex "\def\showsolutions{}% =====================================================================
%  MSACL DS301 Deep Learning · Quiz 5 · Convolutional Networks
%  Academic problem-set style (single-source, \ifsolution toggle)
%  SINGLE SOURCE — produces BOTH the student quiz and the answer key.
%
%  ANSWER TOGGLE
%  -------------
%  Student version (answers HIDDEN) — the default:
%       pdflatex quiz05_cnns.tex
%  Answer key (answers SHOWN):
%       pdflatex "\def\showsolutions{}% =====================================================================
%  MSACL DS301 Deep Learning · Quiz 5 · Convolutional Networks
%  Academic problem-set style (single-source, \ifsolution toggle)
%  SINGLE SOURCE — produces BOTH the student quiz and the answer key.
%
%  ANSWER TOGGLE
%  -------------
%  Student version (answers HIDDEN) — the default:
%       pdflatex quiz05_cnns.tex
%  Answer key (answers SHOWN):
%       pdflatex "\def\showsolutions{}% =====================================================================
%  MSACL DS301 Deep Learning · Quiz 5 · Convolutional Networks
%  Academic problem-set style (single-source, \ifsolution toggle)
%  SINGLE SOURCE — produces BOTH the student quiz and the answer key.
%
%  ANSWER TOGGLE
%  -------------
%  Student version (answers HIDDEN) — the default:
%       pdflatex quiz05_cnns.tex
%  Answer key (answers SHOWN):
%       pdflatex "\def\showsolutions{}\input{quiz05_cnns.tex}"
%  (or, to flip by hand, change \solutionfalse to \solutiontrue below.)
% =====================================================================
\documentclass[11pt]{article}

\newif\ifsolution
\ifdefined\showsolutions \solutiontrue \else \solutionfalse \fi
% \solutiontrue   % <-- uncomment to hard-code the ANSWER KEY build

\usepackage[letterpaper, margin=0.6in, top=0.7in, bottom=0.5in]{geometry}
\usepackage{amsmath, amssymb}
\usepackage{xcolor}
\usepackage{array}
\usepackage{tikz}
\usepackage{fancyhdr}

\pagestyle{fancy}
\fancyhf{}
\fancyhead[L]{\small\textbf{MSACL DS301 --- Deep Learning}}
\fancyhead[R]{\small\textbf{Quiz 5: Convolutional Networks\ifsolution\ \textmd{(Answer Key)}\fi}}
\fancyfoot[C]{\thepage}
\renewcommand{\headrulewidth}{0.4pt}

\newcommand{\blk}[1]{\underline{\hspace{#1}}}
% Reveal-or-blank: shows the answer in the key, a blank rule on the student sheet.
\newcommand{\rb}[2]{\ifsolution\textbf{#1}\else\blk{#2}\fi}
% Boxed answer (key only) and blank writing space (student only).
\newcommand{\keybox}[1]{\ifsolution\par\vspace{2pt}%
  \fbox{\parbox{0.965\textwidth}{\small\textbf{Answer.}\; #1}}\par\fi}
\newcommand{\work}[1]{\ifsolution\else\vspace{#1}\fi}

% ---- numeric grids (course palette) --------------------------------------
\definecolor{ink}{HTML}{1B2025}
\definecolor{teal}{HTML}{0E7C7B}
\definecolor{amber}{HTML}{E09E2F}
% square, centred, monospace cells for the numeric grids
\newcolumntype{G}{>{\centering\arraybackslash\ttfamily}p{6.5mm}}
\newcommand{\gr}{\renewcommand{\arraystretch}{1.5}}

\setlength{\parindent}{0pt}
\setlength{\parskip}{3pt}
\setlength{\abovedisplayskip}{5pt}
\setlength{\belowdisplayskip}{5pt}

\begin{document}

\begin{center}
{\Large \textbf{Quiz 5 --- Convolutional Networks}}\\[2pt]
{\normalsize Convolution by hand, stride, max pooling, output size, and peak-picking as detection}
\end{center}

\vspace{-2pt}
\begin{center}
\fbox{\parbox{0.92\textwidth}{\small
\textbf{Instructions:}\; Answer all four questions --- the same overlay--multiply--add drill you
did on the slides, new numbers. Intuition first; no calculators needed.%
\ifsolution\else\\[4pt]\textbf{Name:}\ \blk{6cm}\hfill\textbf{Date:}\ \blk{3.5cm}\fi
}}
\end{center}

\vspace{-2pt}\hrule\vspace{5pt}

% =========================== Q1 ===========================
\textbf{1. Convolution by hand --- stride 1 and stride 2.}\; {\small \emph{The rule:} lay the
$3\times3$ filter on a $3\times3$ window of the image, multiply each pair of matching cells, and add
the nine products --- that number is one feature-map cell. \emph{Stride} is how far the window jumps;
with a $5\times5$ image and a $3\times3$ filter the map is $(5-3)/s+1$ cells per side.}

\vspace{4pt}
\begin{center}
\begin{minipage}[c]{0.40\linewidth}\centering
{\small image \textbf{5$\times$5}}\\[3pt]
\gr\begin{tabular}{|G|G|G|G|G|}\hline
1 & 1 & 0 & 1 & 0 \\\hline
0 & 1 & 1 & 0 & 1 \\\hline
1 & 0 & 1 & 0 & 0 \\\hline
0 & 1 & 0 & 1 & 1 \\\hline
1 & 0 & 1 & 0 & 1 \\\hline
\end{tabular}
\end{minipage}%
\begin{minipage}[c]{0.30\linewidth}\centering
{\small filter \textbf{3$\times$3}}\\[3pt]
{\color{teal}\gr\begin{tabular}{|G|G|G|}\hline
1 & 0 & 1 \\\hline
0 & 1 & 0 \\\hline
1 & 0 & 1 \\\hline
\end{tabular}}\\[3pt]
{\scriptsize (corners + centre)}
\end{minipage}
\end{center}

\vspace{3pt}
\textbf{(a)} At \emph{stride 1}, the output map is $3\times3$. The cell at output position $(r,c)$ uses
the window whose top-left corner sits at image row $r$, column $c$. Compute the \emph{top row}:

\vspace{2pt}
\quad $(1,1) = $ \rb{$4$}{1.6cm} \hfill $(1,2) = $ \rb{$3$}{1.6cm} \hfill
$(1,3) = $ \rb{$1$}{1.6cm} \hspace*{1.2cm}

\work{1.1cm}

\textbf{(b)} At \emph{stride 2} the map is $2\times2$ (windows start at image rows/cols $1$ and $3$).
Compute its \emph{bottom-right} cell (window top-left at image row $3$, column $3$): \rb{$4$}{1.6cm}

\work{1.0cm}

\keybox{Overlay the filter, multiply, add --- with this filter each cell is just the sum of the four
\emph{corners} and the \emph{centre} of the window.\\[3pt]
(a) $(1,1)$: window rows 1--3, cols 1--3 $=\!\begin{smallmatrix}1&1&0\\0&1&1\\1&0&1\end{smallmatrix}$,
corners+centre $=1+0+1+1+1=\mathbf{4}$.\;
$(1,2)$: cols 2--4 $=\!\begin{smallmatrix}1&0&1\\1&1&0\\0&1&0\end{smallmatrix}$,
$=1+1+0+0+1=\mathbf{3}$.\;
$(1,3)$: cols 3--5 $=\!\begin{smallmatrix}0&1&0\\1&0&1\\1&0&0\end{smallmatrix}$,
$=0+0+1+0+0=\mathbf{1}$.\\[3pt]
(b) Bottom-right stride-2 cell: window rows 3--5, cols 3--5
$=\!\begin{smallmatrix}1&0&0\\0&1&1\\1&0&1\end{smallmatrix}$, corners+centre $=1+0+1+1+1=\mathbf{4}$.
\emph{Stride 2 just samples every other window, so its cells are a subset of the stride-1 map's
values.}}

\vspace{2pt}\hrule\vspace{5pt}

% =========================== Q2 ===========================
\textbf{2. Max pooling.}\; {\small Split the $4\times4$ feature map into non-overlapping $2\times2$
windows (stride 2) and keep each window's \emph{largest} value, giving a $2\times2$ pooled map.}

\vspace{4pt}
\begin{center}
\begin{minipage}[c]{0.40\linewidth}\centering
{\small feature map \textbf{4$\times$4}}\\[3pt]
\gr\begin{tabular}{|G|G|G|G|}\hline
1 & 3 & 2 & 0 \\\hline
4 & 2 & 1 & 5 \\\hline
0 & 1 & 6 & 2 \\\hline
2 & 3 & 1 & 4 \\\hline
\end{tabular}
\end{minipage}%
\begin{minipage}[c]{0.06\linewidth}\centering $\rightarrow$ \end{minipage}%
\begin{minipage}[c]{0.30\linewidth}\centering
{\small pooled \textbf{2$\times$2}}\\[3pt]
\renewcommand{\arraystretch}{1.7}%
\begin{tabular}{|>{\centering\arraybackslash\ttfamily\bfseries}p{9mm}|>{\centering\arraybackslash\ttfamily\bfseries}p{9mm}|}\hline
\rb{4}{6mm} & \rb{5}{6mm} \\\hline
\rb{3}{6mm} & \rb{6}{6mm} \\\hline
\end{tabular}
\end{minipage}
\end{center}

\work{0.5cm}

\keybox{Top-left window $\{1,3,4,2\}\rightarrow\mathbf{4}$;\; top-right $\{2,0,1,5\}\rightarrow\mathbf{5}$;\;
bottom-left $\{0,1,2,3\}\rightarrow\mathbf{3}$;\; bottom-right $\{6,2,1,4\}\rightarrow\mathbf{6}$. So the
pooled map is $\left(\begin{smallmatrix}4&5\\3&6\end{smallmatrix}\right)$. Pooling has no weights ---
it just keeps each region's strongest response, halving the size while a peak that wobbles one cell
over still gives the same answer.}

\vspace{2pt}\hrule\vspace{5pt}

% =========================== Q3 ===========================
\textbf{3. Output size, with padding.}\; {\small A convolution layer's output size is
$\text{out}=(n-k+2p)/s+1$, where $n$ = input size, $k$ = filter size, $p$ = zero-padding per side,
$s$ = stride. A $3\times3$ filter ($k=3$) slides over a length-$7$ input ($n=7$) at stride $s=1$.}

\vspace{3pt}
\textbf{(a)} With padding $p=1$, the output length is\quad
$(7-3+2\cdot1)/1+1=$ \rb{$7$}{2cm}

\work{0.4cm}

\textbf{(b)} With \emph{no} padding ($p=0$), the output length is\quad
$(7-3+0)/1+1=$ \rb{$5$}{2cm}

\work{0.6cm}

\keybox{(a) $(7-3+2)/1+1 = 6+1 = \mathbf{7}$ --- padding~1 with a $3\times3$ filter keeps the length
unchanged (``same'' convolution). (b) $(7-3+0)/1+1 = 4+1 = \mathbf{5}$ --- without padding the map
\emph{shrinks} by two, and the edge positions get fewer looks. This is why we pad: to preserve length
and keep the edges in play.}

\vspace{2pt}\hrule\vspace{5pt}

% =========================== Q4 ===========================
\textbf{4. Peak picking as object detection.}\; {\small A CNN detector run over a chromatogram draws a
\emph{bounding box} with a \emph{confidence score} around each object it finds --- exactly as it would
around a cell in a photo.}

\vspace{3pt}
\textbf{(a)} For a peak on a chromatogram, what does the bounding \emph{box} correspond to?

\work{0.9cm}
\ifsolution\else\blk{\linewidth}\fi

\textbf{(b)} What does the confidence \emph{score} mean?

\work{0.9cm}
\ifsolution\else\blk{\linewidth}\fi

\keybox{(a) The \textbf{box} is the peak's \textbf{retention-time (or $m/z$) window} --- the start-to-end
integration bounds that mark where the peak lies along the axis (its ``where''). (b) The \textbf{score}
is the model's \textbf{confidence (0--1) that the boxed region is a real peak} rather than baseline
noise or an artifact (its ``how sure''). So peak review --- deciding \emph{where} a peak is and
\emph{whether} it is genuine --- is exactly the detection problem (boxes + scores) a CNN already
solves; PeakOnly (2020) does this for real.}

\vspace{4pt}\hrule\vspace{3pt}
{\small\itshape \ifsolution Answer key.\else Keep this sheet --- the convolution, pooling, and
output-size drills are your Lab 2 cheat-sheet.\fi}

\end{document}
"
%  (or, to flip by hand, change \solutionfalse to \solutiontrue below.)
% =====================================================================
\documentclass[11pt]{article}

\newif\ifsolution
\ifdefined\showsolutions \solutiontrue \else \solutionfalse \fi
% \solutiontrue   % <-- uncomment to hard-code the ANSWER KEY build

\usepackage[letterpaper, margin=0.6in, top=0.7in, bottom=0.5in]{geometry}
\usepackage{amsmath, amssymb}
\usepackage{xcolor}
\usepackage{array}
\usepackage{tikz}
\usepackage{fancyhdr}

\pagestyle{fancy}
\fancyhf{}
\fancyhead[L]{\small\textbf{MSACL DS301 --- Deep Learning}}
\fancyhead[R]{\small\textbf{Quiz 5: Convolutional Networks\ifsolution\ \textmd{(Answer Key)}\fi}}
\fancyfoot[C]{\thepage}
\renewcommand{\headrulewidth}{0.4pt}

\newcommand{\blk}[1]{\underline{\hspace{#1}}}
% Reveal-or-blank: shows the answer in the key, a blank rule on the student sheet.
\newcommand{\rb}[2]{\ifsolution\textbf{#1}\else\blk{#2}\fi}
% Boxed answer (key only) and blank writing space (student only).
\newcommand{\keybox}[1]{\ifsolution\par\vspace{2pt}%
  \fbox{\parbox{0.965\textwidth}{\small\textbf{Answer.}\; #1}}\par\fi}
\newcommand{\work}[1]{\ifsolution\else\vspace{#1}\fi}

% ---- numeric grids (course palette) --------------------------------------
\definecolor{ink}{HTML}{1B2025}
\definecolor{teal}{HTML}{0E7C7B}
\definecolor{amber}{HTML}{E09E2F}
% square, centred, monospace cells for the numeric grids
\newcolumntype{G}{>{\centering\arraybackslash\ttfamily}p{6.5mm}}
\newcommand{\gr}{\renewcommand{\arraystretch}{1.5}}

\setlength{\parindent}{0pt}
\setlength{\parskip}{3pt}
\setlength{\abovedisplayskip}{5pt}
\setlength{\belowdisplayskip}{5pt}

\begin{document}

\begin{center}
{\Large \textbf{Quiz 5 --- Convolutional Networks}}\\[2pt]
{\normalsize Convolution by hand, stride, max pooling, output size, and peak-picking as detection}
\end{center}

\vspace{-2pt}
\begin{center}
\fbox{\parbox{0.92\textwidth}{\small
\textbf{Instructions:}\; Answer all four questions --- the same overlay--multiply--add drill you
did on the slides, new numbers. Intuition first; no calculators needed.%
\ifsolution\else\\[4pt]\textbf{Name:}\ \blk{6cm}\hfill\textbf{Date:}\ \blk{3.5cm}\fi
}}
\end{center}

\vspace{-2pt}\hrule\vspace{5pt}

% =========================== Q1 ===========================
\textbf{1. Convolution by hand --- stride 1 and stride 2.}\; {\small \emph{The rule:} lay the
$3\times3$ filter on a $3\times3$ window of the image, multiply each pair of matching cells, and add
the nine products --- that number is one feature-map cell. \emph{Stride} is how far the window jumps;
with a $5\times5$ image and a $3\times3$ filter the map is $(5-3)/s+1$ cells per side.}

\vspace{4pt}
\begin{center}
\begin{minipage}[c]{0.40\linewidth}\centering
{\small image \textbf{5$\times$5}}\\[3pt]
\gr\begin{tabular}{|G|G|G|G|G|}\hline
1 & 1 & 0 & 1 & 0 \\\hline
0 & 1 & 1 & 0 & 1 \\\hline
1 & 0 & 1 & 0 & 0 \\\hline
0 & 1 & 0 & 1 & 1 \\\hline
1 & 0 & 1 & 0 & 1 \\\hline
\end{tabular}
\end{minipage}%
\begin{minipage}[c]{0.30\linewidth}\centering
{\small filter \textbf{3$\times$3}}\\[3pt]
{\color{teal}\gr\begin{tabular}{|G|G|G|}\hline
1 & 0 & 1 \\\hline
0 & 1 & 0 \\\hline
1 & 0 & 1 \\\hline
\end{tabular}}\\[3pt]
{\scriptsize (corners + centre)}
\end{minipage}
\end{center}

\vspace{3pt}
\textbf{(a)} At \emph{stride 1}, the output map is $3\times3$. The cell at output position $(r,c)$ uses
the window whose top-left corner sits at image row $r$, column $c$. Compute the \emph{top row}:

\vspace{2pt}
\quad $(1,1) = $ \rb{$4$}{1.6cm} \hfill $(1,2) = $ \rb{$3$}{1.6cm} \hfill
$(1,3) = $ \rb{$1$}{1.6cm} \hspace*{1.2cm}

\work{1.1cm}

\textbf{(b)} At \emph{stride 2} the map is $2\times2$ (windows start at image rows/cols $1$ and $3$).
Compute its \emph{bottom-right} cell (window top-left at image row $3$, column $3$): \rb{$4$}{1.6cm}

\work{1.0cm}

\keybox{Overlay the filter, multiply, add --- with this filter each cell is just the sum of the four
\emph{corners} and the \emph{centre} of the window.\\[3pt]
(a) $(1,1)$: window rows 1--3, cols 1--3 $=\!\begin{smallmatrix}1&1&0\\0&1&1\\1&0&1\end{smallmatrix}$,
corners+centre $=1+0+1+1+1=\mathbf{4}$.\;
$(1,2)$: cols 2--4 $=\!\begin{smallmatrix}1&0&1\\1&1&0\\0&1&0\end{smallmatrix}$,
$=1+1+0+0+1=\mathbf{3}$.\;
$(1,3)$: cols 3--5 $=\!\begin{smallmatrix}0&1&0\\1&0&1\\1&0&0\end{smallmatrix}$,
$=0+0+1+0+0=\mathbf{1}$.\\[3pt]
(b) Bottom-right stride-2 cell: window rows 3--5, cols 3--5
$=\!\begin{smallmatrix}1&0&0\\0&1&1\\1&0&1\end{smallmatrix}$, corners+centre $=1+0+1+1+1=\mathbf{4}$.
\emph{Stride 2 just samples every other window, so its cells are a subset of the stride-1 map's
values.}}

\vspace{2pt}\hrule\vspace{5pt}

% =========================== Q2 ===========================
\textbf{2. Max pooling.}\; {\small Split the $4\times4$ feature map into non-overlapping $2\times2$
windows (stride 2) and keep each window's \emph{largest} value, giving a $2\times2$ pooled map.}

\vspace{4pt}
\begin{center}
\begin{minipage}[c]{0.40\linewidth}\centering
{\small feature map \textbf{4$\times$4}}\\[3pt]
\gr\begin{tabular}{|G|G|G|G|}\hline
1 & 3 & 2 & 0 \\\hline
4 & 2 & 1 & 5 \\\hline
0 & 1 & 6 & 2 \\\hline
2 & 3 & 1 & 4 \\\hline
\end{tabular}
\end{minipage}%
\begin{minipage}[c]{0.06\linewidth}\centering $\rightarrow$ \end{minipage}%
\begin{minipage}[c]{0.30\linewidth}\centering
{\small pooled \textbf{2$\times$2}}\\[3pt]
\renewcommand{\arraystretch}{1.7}%
\begin{tabular}{|>{\centering\arraybackslash\ttfamily\bfseries}p{9mm}|>{\centering\arraybackslash\ttfamily\bfseries}p{9mm}|}\hline
\rb{4}{6mm} & \rb{5}{6mm} \\\hline
\rb{3}{6mm} & \rb{6}{6mm} \\\hline
\end{tabular}
\end{minipage}
\end{center}

\work{0.5cm}

\keybox{Top-left window $\{1,3,4,2\}\rightarrow\mathbf{4}$;\; top-right $\{2,0,1,5\}\rightarrow\mathbf{5}$;\;
bottom-left $\{0,1,2,3\}\rightarrow\mathbf{3}$;\; bottom-right $\{6,2,1,4\}\rightarrow\mathbf{6}$. So the
pooled map is $\left(\begin{smallmatrix}4&5\\3&6\end{smallmatrix}\right)$. Pooling has no weights ---
it just keeps each region's strongest response, halving the size while a peak that wobbles one cell
over still gives the same answer.}

\vspace{2pt}\hrule\vspace{5pt}

% =========================== Q3 ===========================
\textbf{3. Output size, with padding.}\; {\small A convolution layer's output size is
$\text{out}=(n-k+2p)/s+1$, where $n$ = input size, $k$ = filter size, $p$ = zero-padding per side,
$s$ = stride. A $3\times3$ filter ($k=3$) slides over a length-$7$ input ($n=7$) at stride $s=1$.}

\vspace{3pt}
\textbf{(a)} With padding $p=1$, the output length is\quad
$(7-3+2\cdot1)/1+1=$ \rb{$7$}{2cm}

\work{0.4cm}

\textbf{(b)} With \emph{no} padding ($p=0$), the output length is\quad
$(7-3+0)/1+1=$ \rb{$5$}{2cm}

\work{0.6cm}

\keybox{(a) $(7-3+2)/1+1 = 6+1 = \mathbf{7}$ --- padding~1 with a $3\times3$ filter keeps the length
unchanged (``same'' convolution). (b) $(7-3+0)/1+1 = 4+1 = \mathbf{5}$ --- without padding the map
\emph{shrinks} by two, and the edge positions get fewer looks. This is why we pad: to preserve length
and keep the edges in play.}

\vspace{2pt}\hrule\vspace{5pt}

% =========================== Q4 ===========================
\textbf{4. Peak picking as object detection.}\; {\small A CNN detector run over a chromatogram draws a
\emph{bounding box} with a \emph{confidence score} around each object it finds --- exactly as it would
around a cell in a photo.}

\vspace{3pt}
\textbf{(a)} For a peak on a chromatogram, what does the bounding \emph{box} correspond to?

\work{0.9cm}
\ifsolution\else\blk{\linewidth}\fi

\textbf{(b)} What does the confidence \emph{score} mean?

\work{0.9cm}
\ifsolution\else\blk{\linewidth}\fi

\keybox{(a) The \textbf{box} is the peak's \textbf{retention-time (or $m/z$) window} --- the start-to-end
integration bounds that mark where the peak lies along the axis (its ``where''). (b) The \textbf{score}
is the model's \textbf{confidence (0--1) that the boxed region is a real peak} rather than baseline
noise or an artifact (its ``how sure''). So peak review --- deciding \emph{where} a peak is and
\emph{whether} it is genuine --- is exactly the detection problem (boxes + scores) a CNN already
solves; PeakOnly (2020) does this for real.}

\vspace{4pt}\hrule\vspace{3pt}
{\small\itshape \ifsolution Answer key.\else Keep this sheet --- the convolution, pooling, and
output-size drills are your Lab 2 cheat-sheet.\fi}

\end{document}
"
%  (or, to flip by hand, change \solutionfalse to \solutiontrue below.)
% =====================================================================
\documentclass[11pt]{article}

\newif\ifsolution
\ifdefined\showsolutions \solutiontrue \else \solutionfalse \fi
% \solutiontrue   % <-- uncomment to hard-code the ANSWER KEY build

\usepackage[letterpaper, margin=0.6in, top=0.7in, bottom=0.5in]{geometry}
\usepackage{amsmath, amssymb}
\usepackage{xcolor}
\usepackage{array}
\usepackage{tikz}
\usepackage{fancyhdr}

\pagestyle{fancy}
\fancyhf{}
\fancyhead[L]{\small\textbf{MSACL DS301 --- Deep Learning}}
\fancyhead[R]{\small\textbf{Quiz 5: Convolutional Networks\ifsolution\ \textmd{(Answer Key)}\fi}}
\fancyfoot[C]{\thepage}
\renewcommand{\headrulewidth}{0.4pt}

\newcommand{\blk}[1]{\underline{\hspace{#1}}}
% Reveal-or-blank: shows the answer in the key, a blank rule on the student sheet.
\newcommand{\rb}[2]{\ifsolution\textbf{#1}\else\blk{#2}\fi}
% Boxed answer (key only) and blank writing space (student only).
\newcommand{\keybox}[1]{\ifsolution\par\vspace{2pt}%
  \fbox{\parbox{0.965\textwidth}{\small\textbf{Answer.}\; #1}}\par\fi}
\newcommand{\work}[1]{\ifsolution\else\vspace{#1}\fi}

% ---- numeric grids (course palette) --------------------------------------
\definecolor{ink}{HTML}{1B2025}
\definecolor{teal}{HTML}{0E7C7B}
\definecolor{amber}{HTML}{E09E2F}
% square, centred, monospace cells for the numeric grids
\newcolumntype{G}{>{\centering\arraybackslash\ttfamily}p{6.5mm}}
\newcommand{\gr}{\renewcommand{\arraystretch}{1.5}}

\setlength{\parindent}{0pt}
\setlength{\parskip}{3pt}
\setlength{\abovedisplayskip}{5pt}
\setlength{\belowdisplayskip}{5pt}

\begin{document}

\begin{center}
{\Large \textbf{Quiz 5 --- Convolutional Networks}}\\[2pt]
{\normalsize Convolution by hand, stride, max pooling, output size, and peak-picking as detection}
\end{center}

\vspace{-2pt}
\begin{center}
\fbox{\parbox{0.92\textwidth}{\small
\textbf{Instructions:}\; Answer all four questions --- the same overlay--multiply--add drill you
did on the slides, new numbers. Intuition first; no calculators needed.%
\ifsolution\else\\[4pt]\textbf{Name:}\ \blk{6cm}\hfill\textbf{Date:}\ \blk{3.5cm}\fi
}}
\end{center}

\vspace{-2pt}\hrule\vspace{5pt}

% =========================== Q1 ===========================
\textbf{1. Convolution by hand --- stride 1 and stride 2.}\; {\small \emph{The rule:} lay the
$3\times3$ filter on a $3\times3$ window of the image, multiply each pair of matching cells, and add
the nine products --- that number is one feature-map cell. \emph{Stride} is how far the window jumps;
with a $5\times5$ image and a $3\times3$ filter the map is $(5-3)/s+1$ cells per side.}

\vspace{4pt}
\begin{center}
\begin{minipage}[c]{0.40\linewidth}\centering
{\small image \textbf{5$\times$5}}\\[3pt]
\gr\begin{tabular}{|G|G|G|G|G|}\hline
1 & 1 & 0 & 1 & 0 \\\hline
0 & 1 & 1 & 0 & 1 \\\hline
1 & 0 & 1 & 0 & 0 \\\hline
0 & 1 & 0 & 1 & 1 \\\hline
1 & 0 & 1 & 0 & 1 \\\hline
\end{tabular}
\end{minipage}%
\begin{minipage}[c]{0.30\linewidth}\centering
{\small filter \textbf{3$\times$3}}\\[3pt]
{\color{teal}\gr\begin{tabular}{|G|G|G|}\hline
1 & 0 & 1 \\\hline
0 & 1 & 0 \\\hline
1 & 0 & 1 \\\hline
\end{tabular}}\\[3pt]
{\scriptsize (corners + centre)}
\end{minipage}
\end{center}

\vspace{3pt}
\textbf{(a)} At \emph{stride 1}, the output map is $3\times3$. The cell at output position $(r,c)$ uses
the window whose top-left corner sits at image row $r$, column $c$. Compute the \emph{top row}:

\vspace{2pt}
\quad $(1,1) = $ \rb{$4$}{1.6cm} \hfill $(1,2) = $ \rb{$3$}{1.6cm} \hfill
$(1,3) = $ \rb{$1$}{1.6cm} \hspace*{1.2cm}

\work{1.1cm}

\textbf{(b)} At \emph{stride 2} the map is $2\times2$ (windows start at image rows/cols $1$ and $3$).
Compute its \emph{bottom-right} cell (window top-left at image row $3$, column $3$): \rb{$4$}{1.6cm}

\work{1.0cm}

\keybox{Overlay the filter, multiply, add --- with this filter each cell is just the sum of the four
\emph{corners} and the \emph{centre} of the window.\\[3pt]
(a) $(1,1)$: window rows 1--3, cols 1--3 $=\!\begin{smallmatrix}1&1&0\\0&1&1\\1&0&1\end{smallmatrix}$,
corners+centre $=1+0+1+1+1=\mathbf{4}$.\;
$(1,2)$: cols 2--4 $=\!\begin{smallmatrix}1&0&1\\1&1&0\\0&1&0\end{smallmatrix}$,
$=1+1+0+0+1=\mathbf{3}$.\;
$(1,3)$: cols 3--5 $=\!\begin{smallmatrix}0&1&0\\1&0&1\\1&0&0\end{smallmatrix}$,
$=0+0+1+0+0=\mathbf{1}$.\\[3pt]
(b) Bottom-right stride-2 cell: window rows 3--5, cols 3--5
$=\!\begin{smallmatrix}1&0&0\\0&1&1\\1&0&1\end{smallmatrix}$, corners+centre $=1+0+1+1+1=\mathbf{4}$.
\emph{Stride 2 just samples every other window, so its cells are a subset of the stride-1 map's
values.}}

\vspace{2pt}\hrule\vspace{5pt}

% =========================== Q2 ===========================
\textbf{2. Max pooling.}\; {\small Split the $4\times4$ feature map into non-overlapping $2\times2$
windows (stride 2) and keep each window's \emph{largest} value, giving a $2\times2$ pooled map.}

\vspace{4pt}
\begin{center}
\begin{minipage}[c]{0.40\linewidth}\centering
{\small feature map \textbf{4$\times$4}}\\[3pt]
\gr\begin{tabular}{|G|G|G|G|}\hline
1 & 3 & 2 & 0 \\\hline
4 & 2 & 1 & 5 \\\hline
0 & 1 & 6 & 2 \\\hline
2 & 3 & 1 & 4 \\\hline
\end{tabular}
\end{minipage}%
\begin{minipage}[c]{0.06\linewidth}\centering $\rightarrow$ \end{minipage}%
\begin{minipage}[c]{0.30\linewidth}\centering
{\small pooled \textbf{2$\times$2}}\\[3pt]
\renewcommand{\arraystretch}{1.7}%
\begin{tabular}{|>{\centering\arraybackslash\ttfamily\bfseries}p{9mm}|>{\centering\arraybackslash\ttfamily\bfseries}p{9mm}|}\hline
\rb{4}{6mm} & \rb{5}{6mm} \\\hline
\rb{3}{6mm} & \rb{6}{6mm} \\\hline
\end{tabular}
\end{minipage}
\end{center}

\work{0.5cm}

\keybox{Top-left window $\{1,3,4,2\}\rightarrow\mathbf{4}$;\; top-right $\{2,0,1,5\}\rightarrow\mathbf{5}$;\;
bottom-left $\{0,1,2,3\}\rightarrow\mathbf{3}$;\; bottom-right $\{6,2,1,4\}\rightarrow\mathbf{6}$. So the
pooled map is $\left(\begin{smallmatrix}4&5\\3&6\end{smallmatrix}\right)$. Pooling has no weights ---
it just keeps each region's strongest response, halving the size while a peak that wobbles one cell
over still gives the same answer.}

\vspace{2pt}\hrule\vspace{5pt}

% =========================== Q3 ===========================
\textbf{3. Output size, with padding.}\; {\small A convolution layer's output size is
$\text{out}=(n-k+2p)/s+1$, where $n$ = input size, $k$ = filter size, $p$ = zero-padding per side,
$s$ = stride. A $3\times3$ filter ($k=3$) slides over a length-$7$ input ($n=7$) at stride $s=1$.}

\vspace{3pt}
\textbf{(a)} With padding $p=1$, the output length is\quad
$(7-3+2\cdot1)/1+1=$ \rb{$7$}{2cm}

\work{0.4cm}

\textbf{(b)} With \emph{no} padding ($p=0$), the output length is\quad
$(7-3+0)/1+1=$ \rb{$5$}{2cm}

\work{0.6cm}

\keybox{(a) $(7-3+2)/1+1 = 6+1 = \mathbf{7}$ --- padding~1 with a $3\times3$ filter keeps the length
unchanged (``same'' convolution). (b) $(7-3+0)/1+1 = 4+1 = \mathbf{5}$ --- without padding the map
\emph{shrinks} by two, and the edge positions get fewer looks. This is why we pad: to preserve length
and keep the edges in play.}

\vspace{2pt}\hrule\vspace{5pt}

% =========================== Q4 ===========================
\textbf{4. Peak picking as object detection.}\; {\small A CNN detector run over a chromatogram draws a
\emph{bounding box} with a \emph{confidence score} around each object it finds --- exactly as it would
around a cell in a photo.}

\vspace{3pt}
\textbf{(a)} For a peak on a chromatogram, what does the bounding \emph{box} correspond to?

\work{0.9cm}
\ifsolution\else\blk{\linewidth}\fi

\textbf{(b)} What does the confidence \emph{score} mean?

\work{0.9cm}
\ifsolution\else\blk{\linewidth}\fi

\keybox{(a) The \textbf{box} is the peak's \textbf{retention-time (or $m/z$) window} --- the start-to-end
integration bounds that mark where the peak lies along the axis (its ``where''). (b) The \textbf{score}
is the model's \textbf{confidence (0--1) that the boxed region is a real peak} rather than baseline
noise or an artifact (its ``how sure''). So peak review --- deciding \emph{where} a peak is and
\emph{whether} it is genuine --- is exactly the detection problem (boxes + scores) a CNN already
solves; PeakOnly (2020) does this for real.}

\vspace{4pt}\hrule\vspace{3pt}
{\small\itshape \ifsolution Answer key.\else Keep this sheet --- the convolution, pooling, and
output-size drills are your Lab 2 cheat-sheet.\fi}

\end{document}
"
%  (or, to flip by hand, change \solutionfalse to \solutiontrue below.)
% =====================================================================
\documentclass[11pt]{article}

\newif\ifsolution
\ifdefined\showsolutions \solutiontrue \else \solutionfalse \fi
% \solutiontrue   % <-- uncomment to hard-code the ANSWER KEY build

\usepackage[letterpaper, margin=0.6in, top=0.7in, bottom=0.5in]{geometry}
\usepackage{amsmath, amssymb}
\usepackage{xcolor}
\usepackage{array}
\usepackage{tikz}
\usepackage{fancyhdr}

\pagestyle{fancy}
\fancyhf{}
\fancyhead[L]{\small\textbf{MSACL DS301 --- Deep Learning}}
\fancyhead[R]{\small\textbf{Quiz 5: Convolutional Networks\ifsolution\ \textmd{(Answer Key)}\fi}}
\fancyfoot[C]{\thepage}
\renewcommand{\headrulewidth}{0.4pt}

\newcommand{\blk}[1]{\underline{\hspace{#1}}}
% Reveal-or-blank: shows the answer in the key, a blank rule on the student sheet.
\newcommand{\rb}[2]{\ifsolution\textbf{#1}\else\blk{#2}\fi}
% Boxed answer (key only) and blank writing space (student only).
\newcommand{\keybox}[1]{\ifsolution\par\vspace{2pt}%
  \fbox{\parbox{0.965\textwidth}{\small\textbf{Answer.}\; #1}}\par\fi}
\newcommand{\work}[1]{\ifsolution\else\vspace{#1}\fi}

% ---- numeric grids (course palette) --------------------------------------
\definecolor{ink}{HTML}{1B2025}
\definecolor{teal}{HTML}{0E7C7B}
\definecolor{amber}{HTML}{E09E2F}
% square, centred, monospace cells for the numeric grids
\newcolumntype{G}{>{\centering\arraybackslash\ttfamily}p{6.5mm}}
\newcommand{\gr}{\renewcommand{\arraystretch}{1.5}}

\setlength{\parindent}{0pt}
\setlength{\parskip}{3pt}
\setlength{\abovedisplayskip}{5pt}
\setlength{\belowdisplayskip}{5pt}

\begin{document}

\begin{center}
{\Large \textbf{Quiz 5 --- Convolutional Networks}}\\[2pt]
{\normalsize Convolution by hand, stride, max pooling, output size, and peak-picking as detection}
\end{center}

\vspace{-2pt}
\begin{center}
\fbox{\parbox{0.92\textwidth}{\small
\textbf{Instructions:}\; Answer all four questions --- the same overlay--multiply--add drill you
did on the slides, new numbers. Intuition first; no calculators needed.%
\ifsolution\else\\[4pt]\textbf{Name:}\ \blk{6cm}\hfill\textbf{Date:}\ \blk{3.5cm}\fi
}}
\end{center}

\vspace{-2pt}\hrule\vspace{5pt}

% =========================== Q1 ===========================
\textbf{1. Convolution by hand --- stride 1 and stride 2.}\; {\small \emph{The rule:} lay the
$3\times3$ filter on a $3\times3$ window of the image, multiply each pair of matching cells, and add
the nine products --- that number is one feature-map cell. \emph{Stride} is how far the window jumps;
with a $5\times5$ image and a $3\times3$ filter the map is $(5-3)/s+1$ cells per side.}

\vspace{4pt}
\begin{center}
\begin{minipage}[c]{0.40\linewidth}\centering
{\small image \textbf{5$\times$5}}\\[3pt]
\gr\begin{tabular}{|G|G|G|G|G|}\hline
1 & 1 & 0 & 1 & 0 \\\hline
0 & 1 & 1 & 0 & 1 \\\hline
1 & 0 & 1 & 0 & 0 \\\hline
0 & 1 & 0 & 1 & 1 \\\hline
1 & 0 & 1 & 0 & 1 \\\hline
\end{tabular}
\end{minipage}%
\begin{minipage}[c]{0.30\linewidth}\centering
{\small filter \textbf{3$\times$3}}\\[3pt]
{\color{teal}\gr\begin{tabular}{|G|G|G|}\hline
1 & 0 & 1 \\\hline
0 & 1 & 0 \\\hline
1 & 0 & 1 \\\hline
\end{tabular}}\\[3pt]
{\scriptsize (corners + centre)}
\end{minipage}
\end{center}

\vspace{3pt}
\textbf{(a)} At \emph{stride 1}, the output map is $3\times3$. The cell at output position $(r,c)$ uses
the window whose top-left corner sits at image row $r$, column $c$. Compute the \emph{top row}:

\vspace{2pt}
\quad $(1,1) = $ \rb{$4$}{1.6cm} \hfill $(1,2) = $ \rb{$3$}{1.6cm} \hfill
$(1,3) = $ \rb{$1$}{1.6cm} \hspace*{1.2cm}

\work{1.1cm}

\textbf{(b)} At \emph{stride 2} the map is $2\times2$ (windows start at image rows/cols $1$ and $3$).
Compute its \emph{bottom-right} cell (window top-left at image row $3$, column $3$): \rb{$4$}{1.6cm}

\work{1.0cm}

\keybox{Overlay the filter, multiply, add --- with this filter each cell is just the sum of the four
\emph{corners} and the \emph{centre} of the window.\\[3pt]
(a) $(1,1)$: window rows 1--3, cols 1--3 $=\!\begin{smallmatrix}1&1&0\\0&1&1\\1&0&1\end{smallmatrix}$,
corners+centre $=1+0+1+1+1=\mathbf{4}$.\;
$(1,2)$: cols 2--4 $=\!\begin{smallmatrix}1&0&1\\1&1&0\\0&1&0\end{smallmatrix}$,
$=1+1+0+0+1=\mathbf{3}$.\;
$(1,3)$: cols 3--5 $=\!\begin{smallmatrix}0&1&0\\1&0&1\\1&0&0\end{smallmatrix}$,
$=0+0+1+0+0=\mathbf{1}$.\\[3pt]
(b) Bottom-right stride-2 cell: window rows 3--5, cols 3--5
$=\!\begin{smallmatrix}1&0&0\\0&1&1\\1&0&1\end{smallmatrix}$, corners+centre $=1+0+1+1+1=\mathbf{4}$.
\emph{Stride 2 just samples every other window, so its cells are a subset of the stride-1 map's
values.}}

\vspace{2pt}\hrule\vspace{5pt}

% =========================== Q2 ===========================
\textbf{2. Max pooling.}\; {\small Split the $4\times4$ feature map into non-overlapping $2\times2$
windows (stride 2) and keep each window's \emph{largest} value, giving a $2\times2$ pooled map.}

\vspace{4pt}
\begin{center}
\begin{minipage}[c]{0.40\linewidth}\centering
{\small feature map \textbf{4$\times$4}}\\[3pt]
\gr\begin{tabular}{|G|G|G|G|}\hline
1 & 3 & 2 & 0 \\\hline
4 & 2 & 1 & 5 \\\hline
0 & 1 & 6 & 2 \\\hline
2 & 3 & 1 & 4 \\\hline
\end{tabular}
\end{minipage}%
\begin{minipage}[c]{0.06\linewidth}\centering $\rightarrow$ \end{minipage}%
\begin{minipage}[c]{0.30\linewidth}\centering
{\small pooled \textbf{2$\times$2}}\\[3pt]
\renewcommand{\arraystretch}{1.7}%
\begin{tabular}{|>{\centering\arraybackslash\ttfamily\bfseries}p{9mm}|>{\centering\arraybackslash\ttfamily\bfseries}p{9mm}|}\hline
\rb{4}{6mm} & \rb{5}{6mm} \\\hline
\rb{3}{6mm} & \rb{6}{6mm} \\\hline
\end{tabular}
\end{minipage}
\end{center}

\work{0.5cm}

\keybox{Top-left window $\{1,3,4,2\}\rightarrow\mathbf{4}$;\; top-right $\{2,0,1,5\}\rightarrow\mathbf{5}$;\;
bottom-left $\{0,1,2,3\}\rightarrow\mathbf{3}$;\; bottom-right $\{6,2,1,4\}\rightarrow\mathbf{6}$. So the
pooled map is $\left(\begin{smallmatrix}4&5\\3&6\end{smallmatrix}\right)$. Pooling has no weights ---
it just keeps each region's strongest response, halving the size while a peak that wobbles one cell
over still gives the same answer.}

\vspace{2pt}\hrule\vspace{5pt}

% =========================== Q3 ===========================
\textbf{3. Output size, with padding.}\; {\small A convolution layer's output size is
$\text{out}=(n-k+2p)/s+1$, where $n$ = input size, $k$ = filter size, $p$ = zero-padding per side,
$s$ = stride. A $3\times3$ filter ($k=3$) slides over a length-$7$ input ($n=7$) at stride $s=1$.}

\vspace{3pt}
\textbf{(a)} With padding $p=1$, the output length is\quad
$(7-3+2\cdot1)/1+1=$ \rb{$7$}{2cm}

\work{0.4cm}

\textbf{(b)} With \emph{no} padding ($p=0$), the output length is\quad
$(7-3+0)/1+1=$ \rb{$5$}{2cm}

\work{0.6cm}

\keybox{(a) $(7-3+2)/1+1 = 6+1 = \mathbf{7}$ --- padding~1 with a $3\times3$ filter keeps the length
unchanged (``same'' convolution). (b) $(7-3+0)/1+1 = 4+1 = \mathbf{5}$ --- without padding the map
\emph{shrinks} by two, and the edge positions get fewer looks. This is why we pad: to preserve length
and keep the edges in play.}

\vspace{2pt}\hrule\vspace{5pt}

% =========================== Q4 ===========================
\textbf{4. Peak picking as object detection.}\; {\small A CNN detector run over a chromatogram draws a
\emph{bounding box} with a \emph{confidence score} around each object it finds --- exactly as it would
around a cell in a photo.}

\vspace{3pt}
\textbf{(a)} For a peak on a chromatogram, what does the bounding \emph{box} correspond to?

\work{0.9cm}
\ifsolution\else\blk{\linewidth}\fi

\textbf{(b)} What does the confidence \emph{score} mean?

\work{0.9cm}
\ifsolution\else\blk{\linewidth}\fi

\keybox{(a) The \textbf{box} is the peak's \textbf{retention-time (or $m/z$) window} --- the start-to-end
integration bounds that mark where the peak lies along the axis (its ``where''). (b) The \textbf{score}
is the model's \textbf{confidence (0--1) that the boxed region is a real peak} rather than baseline
noise or an artifact (its ``how sure''). So peak review --- deciding \emph{where} a peak is and
\emph{whether} it is genuine --- is exactly the detection problem (boxes + scores) a CNN already
solves; PeakOnly (2020) does this for real.}

\vspace{4pt}\hrule\vspace{3pt}
{\small\itshape \ifsolution Answer key.\else Keep this sheet --- the convolution, pooling, and
output-size drills are your Lab 2 cheat-sheet.\fi}

\end{document}
